\section{PULSUS Implementation}
\label{sec:software}

The formulation developed in Sec.~\ref{sec:theory} is implemented in
PULSUS, a Python package for direct frequency-domain spectroscopy of
open quantum systems.
The software separates the specification of the molecular system,
laser pulses, and requested spectral coordinates from the numerical
evaluation of the response.
Both the complete light--matter interaction and the controlled
approximations introduced above are exposed through a common
Liouville-space representation.

\subsection{System and Pulse Representation}
\label{sec:software_system}

A molecular model is represented by the \texttt{SpectroscopySystem}
object.
The required inputs are the system Hamiltonian $H$ and the complete
dipole operator $\mu$, together with optional Lindblad collapse
operators and an initial density matrix.
From these quantities, PULSUS constructs the field-free Liouvillian
$\mathcal{L}$ and the complete interaction superoperator $V$ using the
column-stacking convention of Sec.~\ref{sec:theory}.
If no initial state is stored with the system, one may instead be
supplied when a response is evaluated; the stationary state of
$\mathcal{L}$ can also be obtained directly from the system object.

Calculations employing the molecular rotating-wave approximation
additionally require the positive- and negative-frequency components
of the dipole, $\mu_{+}$ and $\mu_{-}$.
These components are supplied explicitly by the user and are converted
to the corresponding interaction superoperators $V_{+}$ and $V_{-}$.
Full-interaction calculations require only the complete dipole
operator and therefore do not require this decomposition.

Each finite-duration laser pulse is represented by a
\texttt{GaussianPulse} object specified by its carrier frequency
$\omega_L$, temporal width $\sigma$, peak field amplitude $E_0$, and
carrier phase $\Phi$.
These parameters determine both the scalar Gaussian spectra used in
the short-pulse model and the Liouvillian dressing operators
$F_j(\mathcal{L},\omega)$ used in the finite-pulse calculations.

\subsection{Direct Cartesian Spectral Evaluation}
\label{sec:software_cartesian}

For a Cartesian spectral grid, PULSUS exploits the separable dependence
of the third-order response on the excitation and detection frequencies.
Let $\{\omega_{1,i}\}_{i=1}^{N_1}$ and
$\{\omega_{3,j}\}_{j=1}^{N_3}$ denote the requested frequency axes.
For the finite-pulse full-interaction model, the excitation-frequency
dependent vectors are constructed as
%
\begin{equation}
|R_i\rangle\rangle
=
F_2(\mathcal{L},\omega_{1,i})
V
G(\omega_{1,i})
V
F_1(\mathcal{L},\omega_{1,i})
|\rho_0\rangle\rangle ,
\label{eq:software_right}
\end{equation}
%
while the detection-frequency dependent covectors are
%
\begin{equation}
\langle\langle L_j|
=
\langle\langle\mu|
G(\omega_{3,j})
V
e^{\mathcal{L}T}
F_3(\mathcal{L},\omega_{3,j}) .
\label{eq:software_left}
\end{equation}
%
The complete spectrum is then obtained from the contractions
%
\begin{equation}
S_{ji}^{(3)}
=
\langle\langle L_j|R_i\rangle\rangle .
\label{eq:software_factorization}
\end{equation}

This factorization avoids independently reconstructing the complete
Liouville-space pathway at every point of the two-dimensional grid.
For a Liouville-space dimension $n$, PULSUS stores the right-hand
objects as an $n\times N_1$ matrix and the left-hand objects as an
$N_3\times n$ matrix.
The full $N_3\times N_1$ spectrum is therefore produced by a final
matrix multiplication.
The waiting-time propagator $e^{\mathcal{L}T}$ is also constructed only
once for a given waiting time.

The frequency-domain resolvents are applied through linear solves rather
than by explicitly forming the inverse matrix.
For each requested frequency, PULSUS solves
%
\begin{equation}
\left[
(\eta-i\omega)I-\mathcal{L}
\right]
|x\rangle\rangle
=
|b\rangle\rangle ,
\label{eq:software_resolvent_solve}
\end{equation}
%
so that $|x\rangle\rangle=G_{\eta}(\omega)|b\rangle\rangle$.
On the detection side, the corresponding transposed linear system is
solved to construct the action of
$\langle\langle\mu|G_{\eta}(\omega_3)$ without explicitly evaluating
the full resolvent matrix.

The same grid-level organization is used for the finite-pulse RWA,
short-pulse RWA, impulsive RWA, and full-interaction impulsive models.
Consequently, the approximation hierarchy of
Sec.~\ref{sec:theory} can be evaluated on a common set of spectral
coordinates without changing the underlying grid representation.

\subsection{Arbitrary Spectral Coordinates}
\label{sec:software_points}

A complete Cartesian grid is not always the most efficient
representation of a multidimensional spectrum.
For example, spectral weight may occupy only a restricted region of the
frequency plane, or a reconstruction procedure may require an
irregular set of sampling coordinates.
PULSUS therefore also supports direct evaluation at arbitrary paired
spectral coordinates.

Given $K$ requested points
%
\begin{equation}
\left\{
(\omega_{1,k},\omega_{3,k})
\right\}_{k=1}^{K},
\end{equation}
%
the paired-coordinate interface evaluates only
%
\begin{equation}
S_k^{(3)}
=
\langle\langle
L(\omega_{3,k})
|
R(\omega_{1,k})
\rangle\rangle ,
\qquad
k=1,\ldots,K,
\label{eq:software_points}
\end{equation}
%
rather than constructing the Cartesian product of all distinct
$\omega_1$ and $\omega_3$ values.

The implementation further identifies repeated values along each
frequency coordinate.
If the requested set contains $N_1^{\mathrm{u}}$ unique excitation
frequencies and $N_3^{\mathrm{u}}$ unique detection frequencies, the
corresponding right- and left-hand Liouville-space objects are
constructed only
$N_1^{\mathrm{u}}$ and $N_3^{\mathrm{u}}$ times, respectively.
The previously constructed objects are then selected and contracted
pairwise according to the requested coordinates.
Consequently, the calculation does not evaluate spectral points that
are absent from the requested set.

This interface makes irregular masks, locally refined regions, and
targeted spectral sampling possible without modifying the underlying
response calculation.
The selection of sampling coordinates is deliberately separated from
the spectroscopy solver: PULSUS evaluates the coordinates supplied by
the user but does not require a particular adaptive-sampling algorithm.
This separation allows different sampling or reconstruction strategies
to be built on the same direct frequency-domain response engine.

\subsection{Field-Sign Sectors and Spectral Conventions}
\label{sec:software_signatures}

PULSUS retains the complete complex-valued third-order response and is
not restricted to the conventional rephasing and nonrephasing sectors.
The standard pathway labels provide the aliases
%
\begin{equation}
\mathrm{NR}\rightarrow(+,-,+),
\qquad
\mathrm{R}\rightarrow(-,+,+),
\end{equation}
%
but an arbitrary field-sign signature
%
\begin{equation}
(s_1,s_2,s_3),
\qquad
s_j\in\{-1,+1\},
\end{equation}
%
may instead be supplied explicitly.
This permits all eight third-order field-sign sectors to be evaluated
within the same finite-pulse or impulsive response framework and is
used below to resolve the non-rotating-wave contributions to the
nonrephasing signal.

For Cartesian calculations, spectra are returned as complex arrays
with the convention
%
\begin{equation}
S_{ji}
=
S^{(3)}(\omega_{3,j},T,\omega_{1,i}),
\end{equation}
%
so that the first array dimension corresponds to the detection
frequency and the second to the excitation frequency.
PULSUS uses signed first-coherence frequencies internally.
Thus, under the convention adopted here, the nonrephasing sector is
evaluated for $\omega_1>0$, whereas the rephasing sector is evaluated
for $\omega_1<0$.
For visualization, the latter may be displayed against the positive
spectroscopic excitation coordinate $-\omega_1$.

The numerical routines preserve the full complex response throughout
the calculation.
Real, imaginary, or magnitude spectra can therefore be selected only
at the analysis or visualization stage rather than being fixed by the
response solver itself.